\documentclass[11pt]{article}
\usepackage[margin=2.4cm]{geometry}
\usepackage{fontspec}
\setmainfont{DejaVu Sans}
\usepackage{juliamono}
\usepackage{fvextra}
\usepackage{xcolor}
\usepackage{xurl}
\usepackage{hyperref}
\usepackage{amssymb}
\newcommand{\checkbox}{$\square$}
\newcommand{\answerchoices}{\checkbox\ YES \quad \checkbox\ NO \quad \checkbox\ CANNOT DETERMINE}
\newcommand{\decisionexplanation}{%
  \par\smallskip\textbf{If NO, give the reason for the mismatch (required). If
  CANNOT DETERMINE, state what context or evidence is missing:}\par
  \rule{\linewidth}{0.4pt}\par\vspace{7pt}\rule{\linewidth}{0.4pt}\par}
\setlength{\parindent}{0pt}\setlength{\parskip}{5pt}
\setlength{\emergencystretch}{2em}
\DefineVerbatimEnvironment{Vsrc}{Verbatim}{fontsize=\small,frame=single,
  breaklines=true,breakanywhere=true}
\DefineVerbatimEnvironment{Vlean}{Verbatim}{fontsize=\scriptsize,frame=single,
  numbers=left,numbersep=5pt,breaklines=true,breakanywhere=true}
\hypersetup{colorlinks=true,urlcolor=blue}

\begin{document}
\begin{center}{\LARGE Proof Verification Document}\\[2pt]{\large arxiv\_1304.1796v2\_\_source\_enriched\_codex}\\[2pt]{\small Quarantined refutations (counterexample already published)}\end{center}
\begin{center}\fcolorbox{blue!60!black}{blue!4}{\parbox{0.92\linewidth}{\textbf{Make three independent decisions.} Follow the evidence chain: original source $\longrightarrow$ challenged claim $\longrightarrow$ proposed counterexample $\longrightarrow$ Lean theorem. A record is valid only when all three links are correct. Complete the human sign-off before reading the automated review appendix; those earlier results are background, not an answer key.}}\end{center}
\textbf{Use only this document, except for the public source link.} Everything needed for Questions 2 and 3 is included here. For Question 1, open the public source link when one is shown. If there is no public link, the complete stored problem and solution are included here. If the information is still not enough, choose CANNOT DETERMINE.
For each question, choose \textbf{YES}, \textbf{NO}, or \textbf{CANNOT DETERMINE}. Choose CANNOT DETERMINE if the document does not provide enough information for you to decide; do not mark YES if you are uncertain. If you find a specific mismatch, choose NO and explain it in the space below that question. Do not rewrite the claim, add assumptions, or substitute a different counterexample.
\textbf{How copied text is shown.} The problem, quoted passage, and proposed counterexample come from the stored record. The Lean theorem statement and Lean files are copied exactly from the certificate. Some long lines continue on the next printed line; this is only page formatting and does not change the copied text. Line numbers are added only to help you find passages. If the stored record contains a character that cannot be printed, the document shows \texttt{[U+XXXX CONTROL]} in its place.
\begin{description}
\item[Challenged claim.] The exact statement from the source that this record says is false.
\item[Proposed counterexample.] A concrete case that should satisfy the claim's assumptions but make its conclusion false.
\item[Lean certificate.] A Lean file proving a formal theorem. Successful compilation proves the theorem as written, not that it is the right translation of the source.
\end{description}
\section*{Question 1 --- Does this proof verification document accurately represent the source?}
\textbf{What to decide.} Compare the challenged passage below with the original source---either the public link or the complete source text in this document. Answer YES only if they make the same claim and include the same assumptions or conditions. If this document changes the meaning, leaves out a condition, or points to a different passage, answer NO.
\subsection*{Original source}
\textbf{Source record:}
\begin{Vsrc}
qedbench_certificate_availability_triage/data/records_arxiv_codex_batch6.jsonl
\end{Vsrc}
\textbf{Public source:}\par
\url{https://arxiv.org/abs/1304.1796v1}\par
\subsection*{Problem or question from the source (verbatim)}
\begin{Vsrc}
Title: The Round Complexity of Small Set Intersection
\end{Vsrc}
\subsection*{Challenged passage recorded from the source (stored text)}
\begin{Vsrc}
& \le p_0 \left(1 + \frac{2}{k - 1}\right)^{a - 1} && \mbox{For sufficiently large $k$} \\
& \le p_0 \left(1 + \frac{2 (a - 1)}{k - 1}\right)
\end{Vsrc}
\subsection*{Surrounding source lines (stored text with added line labels)}
The labels refer to lines in the solution text stored with the source record; the source-line content is not transliterated.
\begin{Vsrc}
0951 | &&\mbox{Using $e^x \geq 1+x$ for all $x$}\\
0952 | &\geq p_0 \left (1+\frac{a-1}{k-1} \right ) \left (1-\frac{1}{k-1} \right)  && \mbox{Using $p_0 = (1 - 1/k)^{k} < e^{-1}$} \\
0953 | & \geq p_0 \left (1+\frac{a-3}{k-1} \right ) &&\mbox{Using $a \le k$}
0954 | \end{align*}
0955 | The proof of the first statement in this part of the lemma follows using $a \geq t$.
0956 | 
0957 | On the other hand:
0958 | \begin{align*}
0959 | \Pr[{\sf OREQ}(f(x),Y) = 0] &  = \left (1- \frac{1}{k-1} \right )^{k-a}\\
0960 | & = \left (1- \frac{1}{k-1} \right )^{k-1} \left(1 - \frac{1}{k - 1}\right)^{- a + 1} \\
0961 | & \le \left( 1 - \frac{1}{k}\right)^k \left(1 - \frac{1}{k - 1}\right)^{- a + 1} \\
0962 | & = p_0 \left(1 - \frac{1}{k - 1}\right)^{- a + 1} \\
0963 | & \le p_0 \left(1 + \frac{2}{k - 1}\right)^{a - 1} && \mbox{For sufficiently large $k$} \\
0964 | & \le p_0 \left(1 + \frac{2 (a - 1)}{k - 1}\right)
0965 | \end{align*}
0966 | 
0967 | \textbf{Part 2.} For the second part we have:
0968 | \begin{align*}
0969 | &\Pr[\agree(f(x),Y) = i] \\
0970 | &= \Pr[\agree(1^a \circ 2^{k - a}, Y) = i] \\
0971 | & = \binom{k - a}{i} (k-1)^{-i} \left(1 - \frac{1}{k - 1}\right)^{k - a - i}  \\
0972 | & \le \binom{k-a}{i} (k -1)^{-i} e^{- \frac{k - a - i}{k-1}} && \mbox{By Proposition~\ref{prop:exp-ineq} with $t = -\frac{k - a - i}{k - 1}$, $n = k - a - i$.} \\
0973 | & = \binom{k-a}{i} e^{-\frac{k - a - i + i \ln (k - 1) (k - 1)}{k-1}} .
0974 | \end{align*}
0975 | Hence we have:
\end{Vsrc}
\textbf{A YES means:} the proof verification document faithfully represents both the source context and the challenged claim, without dropping a condition or changing its meaning.
\section*{Question 2 --- Does the proposed counterexample disprove that claim?}
\textbf{Claim to test (repeated for convenience).}
\begin{Vsrc}
& \le p_0 \left(1 + \frac{2}{k - 1}\right)^{a - 1} && \mbox{For sufficiently large $k$} \\
& \le p_0 \left(1 + \frac{2 (a - 1)}{k - 1}\right)
\end{Vsrc}
\subsection*{How the record says the argument uses this claim}
This is the record's description, not an established fact; check it against the source.
\begin{Vsrc}
The proof asserts that (1+2/(k-1))^(a-1) is at most 1+2(a-1)/(k-1), an inequality used to establish Part 1's upper bound in the Non-Intersection Bias Lemma.
\end{Vsrc}
\subsection*{Proposed counterexample (exact values or construction)}
\begin{Vsrc}
Set a=3 and p0=1. For every integer k>1, letting d=k-1>0 gives (1+2/d)^(a-1)=(1+2/d)^2=1+4/d+4/d^2, whereas 1+2(a-1)/d=1+4/d. Thus the left side is strictly larger by 4/d^2. Concrete instance: k=1001 gives 251001/250000 > 251000/250000.
\end{Vsrc}
\subsection*{Why the record says the counterexample works}
Check this explanation independently; it is evidence to assess, not an answer key.
\begin{Vsrc}
The failure holds for every k>1, so taking k sufficiently large cannot repair the asserted inequality. It is the reverse of the usual Bernoulli inequality for positive increments and integer exponent at least two.
\end{Vsrc}
\subsection*{How to decide}
\begin{enumerate}
\item List every assumption that applies to the challenged claim.
\item Check that the proposed example satisfies every one of those assumptions.
\item Recompute the claimed outcome and the actual outcome using exact reasoning.
\item Confirm that the conflict is with the challenged claim itself, not an unrelated sentence.
\end{enumerate}
\textbf{A YES means:} the example is fully specified, satisfies all relevant assumptions, and makes the challenged claim false. If a necessary value, condition, or calculation is missing, choose CANNOT DETERMINE.
\section*{Question 3 --- Does the final Lean theorem statement verify this same counterexample?}
\textbf{What to decide.} A Lean file can compile while proving the wrong theorem. Compare these three links independently:
\begin{enumerate}
\item The Lean definitions represent the same mathematical objects, assumptions, and operations as the challenged claim.
\item The Lean construction uses the same concrete values or example shown in Question 2.
\item The final theorem statement asserts that the challenged claim fails on that example, not merely a weaker or nearby statement.
\end{enumerate}
\subsection*{Selected final Lean theorem statement (verbatim UTF-8 excerpt)}
This is an exact contiguous slice of the certificate file, from the theorem keyword through the original \texttt{:=} token where its proof begins. It is shown first so the formal target is not buried in the complete file.
\begin{Vlean}
theorem refutes_sufficiently_large_displayed_inequality :
    (forall N : Nat, exists k a t d : Nat,
      sourceHypotheses N k a t d /\
      0 < pNumerator k d /\
      0 < pDenominator k /\
      0 < commonDenominatorAThree k d /\
      clearedAThreeRight k d < clearedAThreeLeft k d /\
      Not (claimedDisplayedStepAtAThree k d)) /\
    terraConcreteWitness :=
\end{Vlean}
\subsection*{Project-local imported Lean files (complete, verbatim UTF-8)}
This section follows the certificate's import commands recursively and prints every repository-owned Lean source file on which the final theorem statement depends. Each file is copied in full, including its definitions and controls. External Mathlib, Std, and Lean Init source is not duplicated.
No project-local imported Lean files were found. External Mathlib, Std, and Lean Init source is not duplicated in this proof verification document.
\subsection*{Complete counterexample Lean file (verbatim UTF-8)}
The complete Lean file is included below so you can check every definition and supporting result used by the final theorem statement. Use it together with the project-local imported Lean files above.
\textbf{Lean certificate file:}
\begin{Vsrc}
partB_feasibility/arxiv_certificates/arxiv_1304.1796v2__source_enriched_codex.lean
\end{Vsrc}
\begin{Vlean}
/-!
Record: `arxiv_1304.1796v2__source_enriched_codex`

Public source: https://arxiv.org/abs/1304.1796

Exact target scope: the displayed algebraic proof step used in Part 1 of the
Non-Intersection Bias Lemma, namely
`p0 * (1 + 2/(k-1))^(a-1) <=
p0 * (1 + 2*(a-1)/(k-1))` for sufficiently large `k`.
This certificate refutes that step, not the lemma's final probabilistic bound.

Concrete Terra-verified witness: `a = t = 3`, `k = 1001`, `d = k-1 = 1000`.
After cancelling the common positive `p0`, the two normalized sides are
`251001/250000` and `251000/250000`. Equivalently,
`1002^2 = 1004004 > 1004000 = 1000*(1000+4)`.

Material modelling judgments: all parameters are natural numbers. For every
threshold `N`, the witness uses `k = N+4`, so `k>N`, `a=t=3`, `1<=a`,
`t<=a`, and `a<k`. Writing `d=k-1`, the source probability is represented
exactly by numerator `pNum=d^k` and denominator `pDen=k^k`, since
`p0=((k-1)/k)^k`. Both are positive. Multiplication by the positive common
denominator `pDen*d^2` converts the source step exactly to
`pNum*(d+2)^2 <= pNum*d*(d+4)`. The left side instead exceeds the right by
`4*pNum`. Terra's `p0=1` is therefore the cancelled normalization, not a
replacement probability model.

The lemma setup can realize `a=3` by taking `x` constantly one and `f(x)`
one on exactly three coordinates and two elsewhere. Neither that realization
nor the conditional law of `Y` is load-bearing for this pointwise algebraic
inference, so no finite probability surrogate is introduced. The negative
control changes the witness feature to `a=2`; after clearing its positive
common denominator, both sides are exactly `pNum*(d+2)` for every `d>0`.
-/

namespace Arxiv13041796v2SourceEnrichedCodex

def witnessA : Nat := 3
def witnessT : Nat := 3
def witnessD (N : Nat) : Nat := N + 3
def witnessK (N : Nat) : Nat := witnessD N + 1

def pNumerator (k d : Nat) : Nat := d ^ k
def pDenominator (k : Nat) : Nat := k ^ k

/- This is the strictly positive denominator used to clear both `p0` and the
square of `d` in the source expression at `a=3`. -/
def commonDenominatorAThree (k d : Nat) : Nat :=
  pDenominator k * d ^ 2

def clearedAThreeLeft (k d : Nat) : Nat :=
  pNumerator k d * (d + 2) ^ 2

def clearedAThreeRight (k d : Nat) : Nat :=
  pNumerator k d * (d * (d + 4))

def claimedDisplayedStepAtAThree (k d : Nat) : Prop :=
  clearedAThreeLeft k d <= clearedAThreeRight k d

def sourceHypotheses
    (N k a t d : Nat) : Prop :=
  N < k /\
  k = N + 4 /\
  a = 3 /\
  t = 3 /\
  1 <= a /\
  t <= a /\
  a < k /\
  d = k - 1 /\
  0 < d

/- The elementary identity is the cleared, unscaled source comparison. -/
theorem square_gap (d : Nat) :
    (d + 2) ^ 2 = d * (d + 4) + 4 := by
  simp [Nat.pow_succ, Nat.mul_add, Nat.add_mul, Nat.add_assoc,
    Nat.add_comm, Nat.add_left_comm]
  rw [Nat.mul_comm 2 d]
  rw [show d * 4 = d * 2 + d * 2 by
    exact Nat.mul_add d 2 2]
  simp [Nat.add_assoc, Nat.add_comm, Nat.add_left_comm]

theorem cleared_gap (k d : Nat) :
    clearedAThreeLeft k d =
      clearedAThreeRight k d + pNumerator k d * 4 := by
  unfold clearedAThreeLeft clearedAThreeRight
  rw [square_gap]
  rw [Nat.mul_add]

theorem strict_reverse_when_d_positive
    (k d : Nat) (hd : 0 < d) :
    clearedAThreeRight k d < clearedAThreeLeft k d := by
  have hp : 0 < pNumerator k d := by
    exact Nat.pow_pos hd
  rw [cleared_gap]
  exact Nat.lt_add_of_pos_right (Nat.mul_pos hp (by decide))

def terraConcreteWitness : Prop :=
  witnessK 997 = 1001 /\
  witnessD 997 = 1000 /\
  1002 ^ 2 = 1004004 /\
  1000 * (1000 + 4) = 1004000 /\
  1004000 < 1004004 /\
  Not (1002 ^ 2 <= 1000 * (1000 + 4)) /\
  1000 ^ 2 = 1000000 /\
  1004004 = 4 * 251001 /\
  1000000 = 4 * 250000 /\
  1004000 = 4 * 251000 /\
  251000 < 251001

/-- Load-bearing refutation: every threshold has a source-admissible
`a=t=3`, `k=N+4` witness. The actual source numerator and denominator and the
clearing denominator are positive, while the exact cleared source inequality
is strictly reversed. The final conjunct records the concrete Terra witness. -/
theorem refutes_sufficiently_large_displayed_inequality :
    (forall N : Nat, exists k a t d : Nat,
      sourceHypotheses N k a t d /\
      0 < pNumerator k d /\
      0 < pDenominator k /\
      0 < commonDenominatorAThree k d /\
      clearedAThreeRight k d < clearedAThreeLeft k d /\
      Not (claimedDisplayedStepAtAThree k d)) /\
    terraConcreteWitness := by
  constructor
  · intro N
    let d := witnessD N
    let k := witnessK N
    have hd : 0 < d := by
      simp [d, witnessD]
    have hk : 0 < k := by
      simp [k, witnessK, witnessD]
    have hpNum : 0 < pNumerator k d := by
      exact Nat.pow_pos hd
    have hpDen : 0 < pDenominator k := by
      exact Nat.pow_pos hk
    have hcommon : 0 < commonDenominatorAThree k d := by
      unfold commonDenominatorAThree
      exact Nat.mul_pos hpDen (Nat.pow_pos hd)
    have hreverse :
        clearedAThreeRight k d < clearedAThreeLeft k d :=
      strict_reverse_when_d_positive k d hd
    refine ⟨k, witnessA, witnessT, d, ?_⟩
    constructor
    · simp [sourceHypotheses, k, d, witnessK, witnessD,
        witnessA, witnessT, Nat.add_assoc]
    constructor
    · exact hpNum
    constructor
    · exact hpDen
    constructor
    · exact hcommon
    constructor
    · exact hreverse
    · unfold claimedDisplayedStepAtAThree
      exact Nat.not_le_of_lt hreverse
  · unfold terraConcreteWitness
    constructor
    · decide
    constructor
    · decide
    constructor
    · decide
    constructor
    · decide
    constructor
    · decide
    constructor
    · decide
    constructor
    · decide
    constructor
    · decide
    constructor
    · decide
    constructor
    · decide
    · decide

def negativeControlA : Nat := 2

def commonDenominatorATwo (k d : Nat) : Nat :=
  pDenominator k * d

def clearedATwoLeft (k d : Nat) : Nat :=
  pNumerator k d * (d + 2)

def clearedATwoRight (k d : Nat) : Nat :=
  pNumerator k d * (d + 2)

/-- Meaningful changed-feature negative control: setting `a=2` changes the
exponent from two to one. For every positive `d`, with `k=d+1`, the source
`p0` numerator, denominator, and clearing denominator remain positive, and
the two exactly cleared sides coincide. -/
theorem negative_control_a_two_equality
    (d : Nat) (hd : 0 < d) :
    negativeControlA = 2 /\
    0 < pNumerator (d + 1) d /\
    0 < pDenominator (d + 1) /\
    0 < commonDenominatorATwo (d + 1) d /\
    clearedATwoLeft (d + 1) d = clearedATwoRight (d + 1) d := by
  have hk : 0 < d + 1 := by
    exact Nat.zero_lt_succ d
  have hpNum : 0 < pNumerator (d + 1) d := by
    exact Nat.pow_pos hd
  have hpDen : 0 < pDenominator (d + 1) := by
    exact Nat.pow_pos hk
  constructor
  · rfl
  constructor
  · exact hpNum
  constructor
  · exact hpDen
  constructor
  · unfold commonDenominatorATwo
    exact Nat.mul_pos hpDen hd
  · rfl

end Arxiv13041796v2SourceEnrichedCodex

\end{Vlean}
\subsection*{Proof-safety information}
The automated scan found no unfinished proofs, new axiom declarations, unsafe declarations, external code hooks, or uses of Lean's compiled decision procedure. This scan concerns proof construction only; it does not establish that the theorem matches the source claim.
\textit{Technical scan counts:} native\_decide=0, sorry=0, axiom\_decl=0, unsafe=0, extern=0. \textit{Technical tier:} KERNEL-ONLY.
\textbf{A YES means:} Lean represents the same objects and assumptions, uses the same proposed counterexample, and the final theorem statement asserts the same contradiction. Successful compilation by itself is not enough.
\clearpage
\section*{Human sign-off}
Complete this block from your own review of Questions 1--3. Do not use the automated appendix below as an answer key.\par
\textbf{Reviewer name:} \rule{9cm}{0.4pt}\par
\textbf{Date:} \rule{4cm}{0.4pt}\par
\medskip\textbf{1. Does the proof verification document accurately represent the source and challenged claim?}\par
\answerchoices\par
\decisionexplanation
\medskip\textbf{2. Does the proposed counterexample satisfy the assumptions and make that claim false?}\par
\answerchoices\par
\decisionexplanation
\medskip\textbf{3. Does the final Lean theorem statement represent and verify that same counterexample?}\par
\answerchoices\par
\decisionexplanation
\medskip\textbf{Overall decision:}\par
\checkbox\ VALIDATED (all three answers are YES)\par
\checkbox\ REJECTED (at least one answer is NO)\par
\checkbox\ NEEDS ANOTHER REVIEWER (at least one CANNOT DETERMINE and no NO)\par
\medskip\textbf{Additional comments.} Use this space for any extra source, counterexample, calculation, or Lean details that support your decisions.\par\rule{\linewidth}{0.4pt}\par\vspace{7pt}\rule{\linewidth}{0.4pt}
\clearpage
\section*{Appendix --- Prior automated review (background only)}
\textbf{These results are not an answer key.} They come from the earlier Opus/Fable verification campaign and are printed only so an auditor can see what was previously checked. The human decisions above must be made independently. Raw status codes are retained in small type for traceability.
\textbf{Overall automated disposition:} Cleared after a follow-up reviewer resolved an automated concern.
\par\textit{Audit code:} \nolinkurl{PASS_ADJUDICATED}
\subsection*{Independent search with no access to the stored example}
\textbf{Result:} An independent search found a counterexample without seeing the stored example.
\par\textit{Audit code:} \nolinkurl{WITNESS_FOUND}
\par\textbf{Checker's explanation:}
\begin{Vsrc}
Setting x = 2/(k-1) > 0 and n = a-1, the quoted step asserts p_0*(1+x)^n <= p_0*(1+n*x). This is Bernoulli's inequality with the direction reversed: for x > 0 and integer n >= 2 the true relation is (1+x)^n > 1 + n*x, since the binomial expansion (1+x)^n = 1 + n*x + C(n,2)*x^2 +... contributes strictly positive higher-order terms. Multiplying by p_0 >= 0 preserves the direction, so the step is false for every k >= 2 whenever a >= 3 and p_0 > 0. (Shortened here; the complete note remains in the machine-check JSON.)
\end{Vsrc}
\par\textbf{Cases checked:}
\begin{Vsrc}
Exhaustive product sweep over k in {2,...,400} U {500, 1000, 2000, 5000, 10000, 20000} (306 values, k >= 2 so that k-1 >= 1 and the expression is defined), a in {1,...,60} (integer round/step count, exponent a-1 >= 0), and p_0 in the 11 exact rationals {0, 1/10^6, 1/1000, 1/100, 1/10, 1/4, 1/3, 1/2, 2/3, 3/4, 1} (p_0 is a probability, so 0 <= p_0 <= 1). (Shortened here; the complete note remains in the machine-check JSON.)
\end{Vsrc}
\subsection*{Independent re-check of the proposed counterexample}
\textbf{Result:} A separate checker recomputed the proposed counterexample and confirmed it.
\par\textit{Audit code:} \nolinkurl{CONFIRMED}
\par\textbf{Checker's note:}
\begin{Vsrc}
Independent recomputation reproduces the recorded witness exactly. With p0=1, a=3, k=1001: LHS = (1+2/1000)^2 = 251001/250000 and RHS = 1+4/1000 = 251000/250000, so LHS - RHS = 1/250000 > 0 and the quoted '<=' is violated. The witness is fully specified and finitary; nothing was missing. Checker exits 0 and prints WITNESS REFUTES CLAIM: True.
\end{Vsrc}
\subsection*{Independent review of the source-to-Lean correspondence}
\textbf{Result:} The Lean proof establishes a related failure, but at least one part of the correspondence needs review.
\par\textit{Audit code:} \nolinkurl{PARTIAL}
\par\textbf{Theorem summary from that reviewer:}
\begin{Vsrc}
The load-bearing theorem `refutes_sufficiently_large_displayed_inequality` is a conjunction of two parts, all over `Nat` with zero imports. Part one: for every threshold `N` there exist `k a t d` with `sourceHypotheses N k a t d` (i.e. `N < k`, `k = N+4`, `a = 3`, `t = 3`, `1 <= a`, `t <= a`, `a < k`, `d = k-1`, `0 < d`), such that `pNumerator k d = d^k > 0`, `pDenominator k = k^k > 0`, `commonDenominatorAThree k d = k^k * d^2 > 0`, and the denominator-cleared source comparison is STRICTLY. (Shortened here; the complete note remains in the machine-check JSON.)
\end{Vsrc}
\par\textbf{Difference noted by that reviewer:}
\begin{Vsrc}
The refutation is REAL, mathematically correct, and the quantifier structure is right (stronger than needed). Verified independently: with p0 = ((k-1)/k)^k, a = 3, k = 1001, the LHS p0*(1+2/1000)^2 normalizes to 251001/250000 and the RHS p0*(1+4/1000) to 251000/250000, so the quoted step fails; multiplying by the positive common denominator k^k * d^2 turns the two sides EXACTLY into pNum*(d+2)^2 and pNum*d*(d+4), confirming the file's clearing algebra with no slack. (Shortened here; the complete note remains in the machine-check JSON.)
\end{Vsrc}
\subsection*{Follow-up review of an automated concern}
\textbf{Result:} Lean uses a documented simplification judged not to change the refutation.
\par\textit{Audit code:} \nolinkurl{DOCUMENTED_MODELING_CHOICE}.\par\textit{Confidence:} \nolinkurl{high}.\par\textit{Recommended action:} \nolinkurl{NONE}.
\par\textbf{Conclusion of the follow-up review:}
\begin{Vsrc}
The quoted step is false. Writing d = k-1 > 0, x = 2/d and n = a-1, it asserts p_0(1+x)^n <= p_0(1+nx), which is Bernoulli's inequality with the direction reversed: for x > 0 and integer n >= 2 the binomial expansion gives (1+x)^n = 1+nx+C(n,2)x^2+... > 1+nx. At a=3 the exact excess is p_0*4/(k-1)^2 > 0 for EVERY integer k > 1, so 'for sufficiently large k' cannot rescue it (and that annotation is in fact attached to the preceding line, not to this one). The record's witness is exact and correct: a=3, p_0=1, k=1001 gives LHS = (1+2/1000)^2 = 251001/250000 and RHS = 1+4/1000 = 251000/250000, gap 1/250000 = 4/1000^2. (Shortened here; the complete note remains in the machine-check JSON.)
\end{Vsrc}
\textit{This appendix is only background. You do not need any separate machine-output files to answer the three questions.}
\end{document}